% from ega2-1-fr.tex
\oldpage[II]{84}

  \item[\rm{b')}] Le morphisme canonique
  $G(\varepsilon)\to\operatorname{Proj}(S)$ est partout défini et est un
  homéomorphisme de l'espace sous-jacent à $X$ sur un sous-espace de
  $\operatorname{Proj}(S)$.
  \item[\rm{c)}] Pour tout $\mathcal{O}_X$-Module quasi-cohérent
  $\mathcal{F}$, si on désigne par $\mathcal{F}_n$ le
  sous-$\mathcal{O}_X$-Module de $\mathcal{F}(n)$ engendré par les sections de
  $\mathcal{F}(n)$ au-dessus de $X$, alors $\mathcal{F}$ est la somme des
  sous-$\mathcal{O}_X$-Modules $\mathcal{F}_n(-n)$, pour les entiers $n>0$.
  \item[\rm{c')}] La propriété c) a lieu pour tout faisceau quasi-cohérent
  d'idéaux de $\mathcal{O}_X$.
\end{enumerate}
En outre, si $(f_\alpha)$ est une famille d'éléments homogènes de $S_+$,
telle que les $X_{f_\alpha}$ soient affines, alors la restriction à
$\bigcup_\alpha X_{f_\alpha}$ du morphisme canonique
$X\to\operatorname{Proj}(S)$ est un isomorphisme de
$\bigcup_\alpha X_{f_\alpha}$ sur
$\bigcup_\alpha(\operatorname{Proj}(S))_{f_\alpha}$.
\end{theorem}

\begin{proof}
Il est clair que b) implique b'), et b') implique a) en vertu de la formule
(\hyperref[II.3.7.3.1-fr]{3.7.3.1}) (tenant compte de ce que
$\varepsilon^\flat$ est l'identité). La condition a) implique a'), car tout
$x\in X$ a un voisinage affine $U$ tel que
$\mathcal{L}|U$ soit isomorphe à $\mathcal{O}_X|U$~; si
$f\in\Gamma(X,\mathcal{L}^{\otimes n})$ est telle que $x\in X_f\subset U$,
$X_f$ est aussi l'ensemble des $x'\in U$ tels que $(f|U)(x')\neq0$, et c'est
par suite un ouvert affine (\hyperref[I.1.3.6-fr]{I, 1.3.6}). Pour prouver
que a') entraîne b), il suffit de prouver la dernière assertion de l'énoncé,
et de vérifier de plus que si $X=\bigcup_\alpha X_{f_\alpha}$, la condition
(iv) de (\hyperref[II.3.8.2-fr]{3.8.2}) est satisfaite. Ce dernier point
résulte aussitôt de (\hyperref[I.9.3.1-fr]{I, 9.3.1, (i)}). Quant à la
dernière assertion de (\hyperref[II.4.5.2-fr]{4.5.2}), comme
$X_{f_\alpha}$ est l'image réciproque de
$(\operatorname{Proj}(S))_{f_\alpha}$ par
$G(\varepsilon)\to\operatorname{Proj}(S)$, il suffit d'appliquer
(\hyperref[I.9.3.2-fr]{I, 9.3.2}). Donc a), a'), b), b') sont équivalentes.

Pour montrer que a') entraîne c), notons que si $X_f$ est affine (avec
$f\in S_k$), $\mathcal{F}|X_f$ est engendré par ses sections au-dessus de
$X_f$ (\hyperref[I.1.3.9-fr]{I, 1.3.9})~; d'autre part
(\hyperref[I.9.3.1-fr]{I, 9.3.1, (ii)}) une telle section $s$ est de la forme
$(t|X_f)\otimes(f|X_f)^{-m}$ où
$t\in\Gamma(X,\mathcal{F}(km))$~; par définition, $t$ est aussi une section
de $\mathcal{F}_{km}$, donc $s$ est bien une section de
$\mathcal{F}_{km}(-km)$ au-dessus de $X_f$, ce qui prouve c). Il est clair
que c) implique c'), et il reste à montrer que c') entraîne a). Or, soit $U$
un voisinage ouvert de $x\in X$, et soit $\mathcal{J}$ un faisceau d'idéaux
quasi-cohérent de $\mathcal{O}_X$ définissant un sous-préschéma fermé de $X$
ayant $X-U$ comme espace sous-jacent
(\hyperref[I.5.2.1-fr]{I, 5.2.1}). L'hypothèse c') entraîne qu'il existe un
entier $n>0$ et une section $f$ de $\mathcal{J}(n)$ au-dessus de $X$ telle
que $f(x)\neq0$. Mais on a évidemment $f\in S_n$, et $x\in X_f\subset U$, ce
qui prouve a).
\end{proof}

Lorsque $X$ est un préschéma dont l'espace sous-jacent est noethérien, les
conditions équivalentes de (\hyperref[II.4.5.2-fr]{4.5.2}) impliquent que
$X$ est un schéma, puisqu'il est isomorphe à un sous-préschéma du schéma
$S=\operatorname{Proj}(A)$ en vertu de
(\hyperref[II.4.5.2-fr]{4.5.2, b)}).

\begin{definition}[4.5.3]
\label{II.4.5.3-fr}
On dit qu'un $\mathcal{O}_X$-Module inversible $\mathcal{L}$ est ample si
$X$ est un schéma quasi-compact et si les conditions équivalentes de
(\hyperref[II.4.5.2-fr]{4.5.2}) sont vérifiées.
\end{definition}

Il résulte évidemment du critère (\hyperref[II.4.5.2-fr]{4.5.2, a)}) que si
$\mathcal{L}$ est un $\mathcal{O}_X$-Module ample, alors, pour tout ouvert
$U$ de $X$, $\mathcal{L}|U$ est un $(\mathcal{O}_X|U)$-Module ample.

Il résulte de la démonstration de (\hyperref[II.4.5.2-fr]{4.5.2}) que les
$X_f$ affines forment déjà une base de la topologie de $X$. De plus~:

\begin{corollary}[4.5.4]
\label{II.4.5.4-fr}
Soit $\mathcal{L}$ un $\mathcal{O}_X$-Module ample. Pour tout sous-espace fini
$Z$ de $X$ et tout voisinage $U$ de $Z$, il existe un entier $n$ et un
$f\in\Gamma(X,\mathcal{L}^{\otimes n})$ tels que $X_f$ soit un voisinage
affine de $Z$ contenu dans $U$.
\end{corollary}

\begin{proof}
